\textbf{Анотація.} Цифрова трансформація наукового адміністрування закладів вищої освіти (ЗВО) стримується фрагментацією наукометричних даних, розпорошених між ORCID, Scopus, Web of Science та OpenAlex. Неузгодженість ідентифікаторів і метаданих призводить до того, що 15--25\,\% записів у типових звітах є дублікатами, а підготовка річного звіту кафедри потребує 8--12 годин ручної праці на копіювання, верифікацію та форматування під вимоги МОН. Управління наукою лишається реактивним: рішення ухвалюються post factum за бібліометричними зведеннями, а науково обґрунтованих механізмів інтеграції генеративного штучного інтелекту (GenAI) в адміністрування досі бракує. Мета статті --- обґрунтувати концептуальну модель і механізми інтеграції GenAI та автономних агентів для переходу від реактивної бібліометричної звітності до проактивного управління науковою діяльністю ЗВО. Запропоновано архітектуру, у якій автономний агент через протокол MCP (Model Context Protocol, Function Calling) звертається до зовнішніх наукометричних джерел, а ядром обробки даних є «золотий запис» (Golden Record) --- конвеєр розв'язання сутностей (Entity Resolution) з канонізацією метаданих, блокуванням, парним логістичним зіставленням, кластеризацією та злиттям атрибутів із провенансом і смугою перевірки людиною в контурі (human-in-the-loop). Модель апробовано на відтворюваному синтетичному бенчмарку, що моделює узагальнений середній український ЗВО (150 науково-педагогічних працівників, 20 кафедр) без реальних персональних даних; статистичну стійкість оцінок забезпечено методом Монте-Карло (K=30 повторень). Конвеєр «золотого запису» досяг показника F1 76,7\,\% проти 59,9\,\% для базового підходу DOI-only та 53,5\,\% для Title+Year. Виправлено критичний дефект нормалізатора, що знищував 100\,\% україномовних назв публікацій. Аналітична модель ефективності оцінює ручну працю в $\approx$11 годин на звіт проти секунд машинного обчислення. Результати позиціонують підхід як реалістичну стратегію цифрової трансформації наукового адміністрування ЗВО засобами ІКТ в управлінні освітою.

\textbf{Ключові слова:} ІКТ в управлінні освітою; цифрова трансформація наукового адміністрування; генеративний штучний інтелект; автономний агент (MCP); золотий запис і розв'язання сутностей; проактивне управління науковою діяльністю ЗВО.

\bigskip

\textbf{Abstract.} The digital transformation of research administration in higher education institutions (HEIs) is hindered by the fragmentation of scientometric data scattered across ORCID, Scopus, Web of Science, and OpenAlex. Inconsistent identifiers and metadata cause 15--25\,\% of records in typical reports to be duplicates, while preparing an annual departmental report demands 8--12 hours of manual labour for copying, verification, and formatting to ministry requirements. Research management remains reactive: decisions are made post factum from bibliometric summaries, and scientifically grounded mechanisms for embedding generative artificial intelligence (GenAI) into administration are still lacking. The aim of the article is to substantiate a conceptual model and integration mechanisms for GenAI and autonomous agents that enable a shift from reactive bibliometric reporting to proactive management of HEI research activity. We propose an architecture in which an autonomous agent queries external scientometric sources through the Model Context Protocol (MCP, Function Calling), while the data-processing core is the Golden Record --- an Entity Resolution pipeline comprising metadata canonicalisation, blocking, pairwise logistic matching, clustering, and provenance-aware attribute merging with a human-in-the-loop review band. The model is evaluated on a reproducible synthetic benchmark that simulates a generalised average Ukrainian HEI (150 academic staff, 20 departments) with no real personal data; statistical robustness is ensured by Monte Carlo simulation (K=30 replications). The Golden Record pipeline attains an F1 score of 76,7\,\% versus 59,9\,\% for the DOI-only baseline and 53,5\,\% for Title+Year. A critical normaliser defect that destroyed 100\,\% of Ukrainian-language publication titles was corrected. An analytical efficiency model estimates manual effort at $\approx$11 hours per report against seconds of machine computation. The findings position the approach as a realistic ICT-in-education-management strategy for the digital transformation of HEI research administration.

\textbf{Keywords:} ICT in education management; digital transformation of research administration; generative artificial intelligence; autonomous agent (MCP); Golden Record and Entity Resolution; proactive management of HEI research activity.